\documentclass[11pt]{article}

\usepackage[margin=1.1in]{geometry}
\usepackage{amsmath,amssymb}
\usepackage{booktabs}
\usepackage{graphicx}
\usepackage{algorithm}
\usepackage{algpseudocode}
\usepackage{xcolor}
\usepackage[colorlinks=true,linkcolor=blue!50!black,citecolor=blue!50!black,urlcolor=blue!50!black]{hyperref}
\usepackage{microtype}

\newcommand{\code}[1]{\texttt{#1}}

\title{\textbf{Blokus v2: A Bitboard Engine and a Budgeted\\ Bandit-Search AI}\\[6pt]
\large Technical Report}
\author{Feng-Ting Liao}
\date{July 2026}

\begin{document}
\maketitle

\begin{abstract}
We describe the engine and artificial intelligence of \emph{Blokus v2}, a Rust
implementation of the four-player board game Blokus. The rules engine encodes
the official $20\times 20$ board as per-player bitboards, reducing legality
testing and full legal-move enumeration to word-parallel bit operations
($\approx$0.2\,ms per move including evaluation). On top of the engine we
build two policies: (i) a fast heuristic policy that scores every legal move
on piece size, corner mobility, opponent-corner denial, and early-game
centrality; and (ii) a strong policy that treats move selection as a budgeted
multi-armed bandit problem, allocating Monte-Carlo rollouts across the top-$K$
heuristic candidates with UCB1 and returning the most-visited arm when a
wall-clock budget expires. In head-to-head evaluation the bandit policy
defeats the heuristic policy in 11 of 12 four-player games while honoring a
one-second decision budget to within a single iteration. The system exposes
the budget to players directly as a ``think-time'' dial ($0$--$5$\,s) in the
game's interface.
\end{abstract}

\section{Introduction}

Blokus is a four-player perfect-information territory game played with
polyomino pieces. Its branching factor is unusually punishing: an opening
position offers each player on the order of $10^3$ legal placements, and
placements interact through both denial (occupying cells opponents need) and
mobility (creating one's own future placement anchors). These properties make
Blokus a pleasant middle ground between classic two-player search domains and
modern multi-agent settings: strong play requires lookahead, but the
four-player, single-move-per-turn structure defeats naive minimax.

This report documents the second major version of the project. Version~1 was
a C++ course project at National Chiao Tung University (now National Yang
Ming Chiao Tung University) by Feng-Ting Liao and ShengHsiung Hsieh; its AI
scanned a $14\times 14$ window around a single frontier cell and placed
near-randomly in the opening. Version~2 is a ground-up Rust rewrite whose AI
is the subject of this report. We cover the board representation
(Section~\ref{sec:engine}), the heuristic evaluation
(Section~\ref{sec:heuristic}), the bandit search
(Section~\ref{sec:bandit}), and empirical results
(Section~\ref{sec:benchmarks}).

\section{Rules and Notation}
\label{sec:rules}

The board is the grid $B=\{0,\dots,19\}^2$. Each player
$p\in\{0,1,2,3\}$ owns the 21 free polyominoes of sizes 1--5 (89 squares
total) and a starting corner $c_p$ (one of the four board corners). Players
move in fixed rotation; a player unable to move is permanently eliminated
from the rotation (opponent pieces can only remove placement options, never
create them, so inability to move is absorbing).

A \emph{placement} is a piece, one of its up to eight orientations
(rotations and reflections, deduplicated per piece), and a translation. Let
$S\subset B$ be the placement's cells and $O_p$ the cells already occupied by
$p$. A placement for $p$ is legal iff
\begin{enumerate}
  \item $S$ lies within $B$ and is disjoint from \emph{all} occupied cells;
  \item $S$ is not 4-adjacent (edge-adjacent) to $O_p$;
  \item $S$ is 8-but-not-4 adjacent (corner-adjacent) to $O_p$; or, if
        $O_p=\varnothing$, $c_p\in S$.
\end{enumerate}
Contact with \emph{other} colors is unrestricted. Terminal scoring is
official: $-1$ per unplaced square, $+15$ if all 21 pieces were placed, and
$+5$ more if the monomino was placed last.

\section{Bitboard Engine}
\label{sec:engine}

\subsection{Representation}

Each player's occupancy $O_p$ is stored as twenty 32-bit words, one per row,
using the low 20 bits ($\approx$80 bytes per player). The two neighborhood
operators used throughout are computed by shifts and ORs over rows:
\begin{align*}
\mathrm{orth}(X) &= X \cup \{\text{cells 4-adjacent to } X\},\\
\mathrm{diag}(X) &= \{\text{cells 8-adjacent but not necessarily 4-adjacent to } X\}
                   \text{ (diagonal spread)}.
\end{align*}
Both are 20-iteration loops of word operations; no per-cell work is done.

\subsection{Forbidden cells and seeds}

For player $p$ define
\begin{align}
F_p &= \Big(\bigcup_q O_q\Big)\;\cup\;\mathrm{orth}(O_p)
      &&\text{(\emph{forbidden}: cells no piece of $p$ may cover)},\\
A_p &= \mathrm{diag}(O_p)\setminus F_p
      &&\text{(\emph{seeds}: legal anchor cells)},
\end{align}
with $A_p=\{c_p\}$ before $p$'s first move. Legality of a placement $S$
reduces to $S\cap F_p=\varnothing$ and $S\cap A_p\neq\varnothing$: rule (1)
and (2) are exactly $S \cap F_p = \varnothing$, and rule (3) is
$S\cap A_p\neq\varnothing$ because any placement satisfying (1)--(2) touches
$O_p$ diagonally iff it covers a seed.

\subsection{Move generation}

Legal moves are enumerated by aligning, for every unused piece, every
orientation cell onto every seed and testing the (at most five) cells of the
placement against $F_p$; duplicates arising from a placement covering several
seeds are removed with a hash set keyed on (piece, orientation, offset). The
count of seeds is small (typically 5--30), so enumeration touches only a tiny
fraction of the $20\times 20\times 8\times 21$ nominal placement space.
$|A_p|$ (\code{seed\_count}) doubles as a cheap mobility measure for
evaluation.

\section{Heuristic Policy}
\label{sec:heuristic}

The fast policy scores every legal move $m$ of player $p$ with
\begin{equation}
\label{eq:eval}
V(m) \;=\; 90\,\mathrm{size}(m)
      \;+\; 14\,\mathrm{corners}(m)
      \;+\; 28\,\mathrm{blocked}(m)
      \;+\; w_c\,\mathrm{center}(m),
\end{equation}
where $\mathrm{size}$ is the piece's square count (points saved under
official scoring); $\mathrm{corners}$ is $|A_p|$ recomputed \emph{after} the
move (future mobility); $\mathrm{blocked}$ is the number of opponent seed
cells covered by the move, summed over active opponents (each is one
placement anchor an opponent loses); and $\mathrm{center}$ sums, over the
move's cells, $9-d_\infty(\text{cell},\text{center }2\times 2)$, with weight
$w_c=4$ during the first five rounds and $1$ afterwards. The weights were
tuned by self-play.

\paragraph{Near-tie sampling.} Instead of a strict $\arg\max$, the policy
samples uniformly among all moves with
$V(m)\ge \max_{m'}V(m')-\tau$, $\tau=25$. This costs no measurable strength
and removes the deterministic ``same opening every game'' artifact (measured
in Section~\ref{sec:diversity}).

\section{Bandit Search}
\label{sec:bandit}

The strong policy treats the decision as a \emph{budgeted best-arm
identification} problem: given a wall-clock budget $T$ (the in-game think
dial, 0--5\,s), spend rollouts where regret arguments say they are most
informative, then commit to one move. Algorithm~\ref{alg:bandit} gives the
full procedure; constants are in Table~\ref{tab:constants}.

\subsection{Arms and priors}

All legal moves are scored with Eq.~\eqref{eq:eval} and the top
$K=40$ become bandit arms. Arm $i$ with static score $s_i$ is initialized
with $n_0=2$ virtual visits at value
\begin{equation}
\hat v_i^{(0)} = 0.35 + 0.3\,\frac{s_i - s_{\min}}{s_{\max}-s_{\min}} + \varepsilon_i,
\qquad \varepsilon_i \sim \mathcal{U}(-0.01, 0.01),
\end{equation}
so that early pulls follow the heuristic ordering without hard-committing to
it; the jitter breaks exact ties among symmetric openings.

\subsection{UCB1 allocation}

At iteration $t$ the search pulls the arm maximizing
\begin{equation}
\mathrm{UCB}_i \;=\; \hat v_i \;+\; C\,\sqrt{\frac{\ln (t+1)}{n_i+1}},
\qquad C = 0.45,
\end{equation}
where $\hat v_i$ is the arm's running mean value in $[0,1]$ and $n_i$ its
visit count (including virtual visits). UCB1's logarithmic-regret guarantee
\cite{auer2002} is what makes the think-time dial behave sensibly: additional
budget is spent almost exclusively on arms still statistically capable of
being best.

\subsection{Rollout policy}

A pull of arm $i$ plays $m_i$, then simulates $H=8$ further plies (two full
rounds) in which \emph{every} player, friend and foe, follows a
\emph{noisy-greedy} policy: it plays
$\arg\max_m \big[V(m) + \sigma G_m\big]$ with i.i.d.\ standard Gumbel noise
$G_m$ and scale $\sigma=25$ evaluation units. By the Gumbel-max
(perturb-and-max) identity this is exactly softmax sampling over $V/\sigma$
\cite{gumbel1954,maddison2014}, giving rollouts principled variety at zero
bookkeeping cost. Plies whose position offers more than 96 legal moves
evaluate a uniform sample of 96 of them, capping rollout latency around
1\,ms with no measured strength loss. Eliminated players are skipped by the
engine's turn logic.

\subsection{Leaf evaluation}

Rollouts are truncated, so leaves are scored by a potential function. For a
terminal leaf the official score differential against the best opponent is
used; otherwise
\begin{equation}
\phi(q) = \mathrm{placed}(q) + 0.4\,|A_q|, \qquad
v \;=\; \tfrac12 + \tfrac12\tanh\!\Big(\frac{\phi(p)-\max_{o\neq p,\,o\,\text{active}}\phi(o)}{12}\Big),
\end{equation}
i.e., squares already banked plus discounted mobility, squashed to $[0,1]$.
Measuring against the \emph{best} opponent rather than the average reflects
that Blokus is won by rank, not by point sum.

\subsection{Final selection}

When the budget expires the search returns the most-visited arm, with one
refinement: arms whose visit count is within 90\% of the maximum are
statistically indistinguishable after a few hundred rollouts, so the policy
picks uniformly among them (exact ties fall back to mean value). This
preserves strength while keeping openings varied between games.

\begin{algorithm}[t]
\caption{Budgeted UCB1 move selection for player $p$}
\label{alg:bandit}
\begin{algorithmic}[1]
\State $M \gets$ \code{legal\_moves}$(p)$;\; \Return the move if $|M|\le 1$; greedy policy if $T=0$
\State score $M$ with Eq.~\eqref{eq:eval}; keep top $K$ as arms with priors $(n_0, \hat v_i^{(0)})$
\While{elapsed $< T$}
  \State $i \gets \arg\max_i\ \hat v_i + C\sqrt{\ln(t+1)/(n_i+1)}$
  \State $G \gets$ copy of root state; play $m_i$
  \For{$h = 1 \dots H$ while $G$ not terminal}
    \State current player plays $\arg\max_m V(m) + \sigma\,\mathrm{Gumbel}()$ \Comment{softmax rollout}
  \EndFor
  \State update arm $i$ with leaf value $v(G)$;\; $t \gets t+1$
\EndWhile
\State \Return uniform choice among arms with $n_i \ge 0.9\max_j n_j$
\end{algorithmic}
\end{algorithm}

\begin{table}[t]
\centering
\caption{Search constants (all tuned by self-play).}
\label{tab:constants}
\begin{tabular}{llr}
\toprule
Constant & Meaning & Value \\
\midrule
$K$ (\code{TOP\_K}) & candidate arms & 40 \\
$C$ (\code{UCB\_C}) & UCB1 exploration & 0.45 \\
$n_0$ (\code{PRIOR\_N}) & virtual prior visits per arm & 2 \\
$H$ (\code{ROLLOUT\_PLIES}) & rollout horizon (plies) & 8 \\
$\sigma$ (\code{GUMBEL\_SCALE}) & rollout softmax temperature & 25 \\
--- (\code{ROLLOUT\_MOVE\_CAP}) & moves evaluated per rollout ply & 96 \\
$\tau$ (\code{GREEDY\_TOLERANCE}) & greedy near-tie band & 25 \\
\bottomrule
\end{tabular}
\end{table}

\section{Benchmarks}
\label{sec:benchmarks}

All measurements: Apple Silicon macOS, \code{cargo run --release} (LTO,
\code{opt-level=3}), single-threaded search. The benchmark driver is
\code{examples/selfplay.rs}; the numbers below are from the run recorded at
the time of writing and are reproduced (within noise) by
\code{cargo run --release --example selfplay}.

\subsection{Engine and heuristic-policy speed}

Twenty full four-player self-play games with the heuristic policy complete in
$\approx$0.25\,s total --- $\approx$\textbf{0.19\,ms per move} including full
legal-move enumeration and evaluation of every candidate. This is the
latency floor of the ``dial = 0'' setting.

\subsection{Heuristic policy vs.\ random baseline}

Seats 0 and 2 play the heuristic policy; seats 1 and 3 play uniformly random
legal moves (comparable to v1's opening behavior). Over 20 games the
heuristic side's best score beats the random side's best score
\textbf{20/20} times, with mean official scores $-8.7$ vs.\ $-37.3$.

\subsection{Bandit search vs.\ heuristic policy}

Seats 0 and 2 play the bandit search at 300 iterations
($\approx$350\,ms per decision); seats 1 and 3 play the heuristic policy.
Over 12 games:

\begin{table}[h]
\centering
\caption{Bandit (300 iterations) vs.\ heuristic policy, 12 four-player games.}
\begin{tabular}{lrr}
\toprule
 & Bandit seats & Heuristic seats \\
\midrule
Wins (best score in game) & \textbf{11/12} & 1/12 \\
Mean official score & $-11.6$ & $-20.0$ \\
\bottomrule
\end{tabular}
\end{table}

\noindent
An 8.4-point mean margin is roughly two pentominoes of extra territory per
game; in one game the bandit seat placed all 21 pieces ($+20$, the maximum
score with monomino-last bonus).

\subsection{Budget compliance}

With \code{Budget::Millis(1000)}, decisions from an opening position and a
round-6 midgame position both returned in 1.001\,s --- the budget is checked
before every rollout, so the worst-case overshoot is a single rollout
($\approx$1--5\,ms). At 300 iterations the mean decision time over 430
decisions was 349\,ms.

\subsection{Opening diversity}
\label{sec:diversity}

Over ten RNG seeds, the number of distinct first moves played as Blue: v2's
original strict-argmax greedy \textbf{1/10} (the user-visible ``the first
pieces are always the same'' defect); near-tie greedy \textbf{8/10}; bandit
search \textbf{5/10}. The bandit's diversity is lower than the greedy's by
design --- rollouts genuinely disambiguate some openings the static
evaluation cannot.

\section{Discussion and Future Work}

The design deliberately spends its budget at the root rather than growing a
tree: with four players and per-ply branching in the hundreds, a
UCT tree \cite{browne2012} of meaningful depth is unreachable in a few
seconds, whereas root-level bandit allocation with strong heuristic rollouts
converts budget into strength immediately and degrades gracefully to the
heuristic policy at $T=0$. Natural extensions, roughly in order of expected
value per effort:

\begin{itemize}
  \item \textbf{Progressive widening / tree growth} below the root for the
        late game, where branching collapses and exact lines matter;
  \item \textbf{Paranoid or max$^n$ backups} in place of the
        best-opponent potential, sharpening play when one opponent is far
        ahead;
  \item \textbf{Learned evaluation} replacing Eq.~\eqref{eq:eval}'s
        hand-tuned weights (the bitboard features are already
        network-friendly);
  \item \textbf{Parallel rollouts} --- the search is single-threaded today;
        arms are independent given the root, so near-linear speedup is
        available.
\end{itemize}

\section{Reproducibility}

\begin{itemize}
  \item Engine and rules tests: \code{cargo test --release} (12 tests: piece
        set, legality, scoring, generator/validator consistency, bandit vs.\
        greedy, budget compliance).
  \item Benchmarks: \code{cargo run --release --example selfplay}.
  \item Interactive play: \code{cargo run --release}; the think-time dial on
        the setup screen sets $T$ per AI move. \code{BLOKUS\_AUTO=<seconds>}
        launches a four-AI game directly.
\end{itemize}

\section*{Acknowledgments}

Version 1 was co-authored with ShengHsiung Hsieh. The v2 Rust translation,
the engine, the search, and this report were produced with Fable~5
(Claude, Anthropic).

\begin{thebibliography}{9}

\bibitem{auer2002}
P.~Auer, N.~Cesa-Bianchi, and P.~Fischer.
\newblock Finite-time analysis of the multiarmed bandit problem.
\newblock \emph{Machine Learning}, 47(2--3):235--256, 2002.

\bibitem{browne2012}
C.~Browne et~al.
\newblock A survey of Monte Carlo tree search methods.
\newblock \emph{IEEE Transactions on Computational Intelligence and AI in
  Games}, 4(1):1--43, 2012.

\bibitem{gumbel1954}
E.~J. Gumbel.
\newblock \emph{Statistical Theory of Extreme Values and Some Practical
  Applications}.
\newblock U.S. National Bureau of Standards, 1954.

\bibitem{maddison2014}
C.~J. Maddison, D.~Tarlow, and T.~Minka.
\newblock A$^*$ sampling.
\newblock In \emph{Advances in Neural Information Processing Systems~27},
  2014.

\bibitem{blokusrules}
Mattel, Inc.
\newblock \emph{Blokus Instruction Sheet R1983}.
\newblock \url{https://service.mattel.com/instruction_sheets/R1983\%20-\%20Blokus.pdf}

\end{thebibliography}

\end{document}
